\begingroup\small
\begin{thebibliography}{99}
\bibitem{furedi}
Z. F\"uredi,
Hypergraphs in which all disjoint pairs have distinct unions,
\emph{Combinatorica} \textbf{4} (1984), 161--168.
\url{https://doi.org/10.1007/BF02579216}.

\bibitem{ff}
P. Frankl and Z. F\"uredi,
A new extremal property of Steiner triple-systems,
\emph{Discrete Mathematics} \textbf{48} (1984), 205--212.

\bibitem{pv}
O. Pikhurko and J. Verstra\"ete,
The maximum size of hypergraphs without generalized 4-cycles,
\emph{Journal of Combinatorial Theory, Series A}
\textbf{116} (2009), 637--649.
\url{https://doi.org/10.1016/j.jcta.2008.09.002}.

\bibitem{mubayi}
D. Mubayi,
Some exact results and new asymptotics for hypergraph Tur\'an numbers,
\emph{Combinatorics, Probability and Computing}
\textbf{11} (2002), 299--309.
\url{https://doi.org/10.1017/S0963548301005028}.
The restricted bound is journal Theorem~1.3; the author manuscript dated
11 June 2005 numbers it Theorem~1.4.

\bibitem{hmy}
Hao Huang, Jie Ma and Tianchi Yang,
Extremal hypergraphs without generalized 4-cycles,
arXiv:2609.37744v1 (29 September 2026).
\url{https://arxiv.org/abs/2609.37744v1}.

\bibitem{howard}
L. Howard,
\emph{Disjoint Union-Free 3-Uniform Hypergraphs},
M.Sc. thesis, University of Victoria, 2006.
\url{https://hdl.handle.net/1828/1868}.

\bibitem{dukesling}
P. J. Dukes and A. C. H. Ling,
Disjoint union-free designs with block size three,
\emph{Bulletin of the Institute of Combinatorics and its Applications}
\textbf{45} (2005), 5--10.

\bibitem{erdos}
Erd\H{o}s Problems, Problem 643,
\url{https://www.erdosproblems.com/643}, accessed 19 October 2026.
\bibitem{erdos1977}
P. Erd\H{o}s, Problems and results in combinatorial analysis,
\emph{Congressus Numerantium} \textbf{19} (1977), 3--12.
\url{https://users.renyi.hu/~p_erdos/1977-24.pdf}.
\bibitem{hanani1961}
H. Hanani, The existence and construction of balanced incomplete block designs,
\emph{Annals of Mathematical Statistics} \textbf{32} (1961), 361--386.
\url{https://doi.org/10.1214/aoms/1177705047}.
The five-block result is (7.10), pp.~378--380, with the 141-point exception completed by the 1965 note.
\bibitem{hanani1965}
H. Hanani, A balanced incomplete block design,
\emph{Annals of Mathematical Statistics} \textbf{36} (1965), 711.
\url{https://doi.org/10.1214/aoms/1177700182}.
This note supplies the 141-point case and removes the exceptions in the 1961 statement (7.10).
\end{thebibliography}
\endgroup
